\documentclass[10pt,a4paper]{amsart}
\usepackage[margin=19mm]{geometry}
\usepackage{amsmath,amssymb,mathtools}
\usepackage[scr=rsfs,cal=euler]{mathalfa}
\usepackage{xfrac}
\usepackage[unicode,hidelinks,bookmarks=false]{hyperref}
\newcommand{\ie}{i.e.\ }
\newcommand{\cf}{cf.\ }
\newcommand{\resp}{respectively\ }
\newcommand{\Subsection}{\subsection}
\providecommand{\nobreakdash}{\nobreak\hbox{-}}
\setlength{\emergencystretch}{2em}
\multlinegap0pt
\usepackage{mathtools}
\usepackage{amssymb}
\usepackage{csquotes}
\usepackage{tikz}
\newtheorem{definition}{Definition}
\newtheorem{lemma}{Lemma}
\newtheorem{theorem}{Theorem}
\usepackage{cleveref}
\crefname{definition}{Definition}{Definitions}
\crefname{lemma}{Lemma}{Lemmas}
\crefname{theorem}{Theorem}{Theorems}
\renewcommand{\geq}{\geqslant}
\renewcommand{\leq}{\leqslant}
\newcommand{\sCop}{\operatorname{sCop}}
\newcommand{\wCop}{\operatorname{wCop}}
\title{Meta gaming in geometric group theory}
\author{Raphael Appenzeller, Kevin Klinge}
\date{Source: arXiv:2502.04540v2 (CC BY 4.0)}
\begin{document}
\maketitle

\noindent\emph{Source and licence.} Meta gaming in geometric group theory. Version arXiv:2502.04540v2 (CC BY 4.0), distributed by arXiv under the Creative Commons Attribution 4.0 International licence, http://creativecommons.org/licenses/by/4.0/, as recorded in the arXiv licence metadata for this exact version and shown on the abstract page https://arxiv.org/abs/2502.04540v2. Full licence notice: \texttt{licenses/arxiv-2502.04540-CC-BY-4.0.txt}.\par\smallskip

\noindent\textit{Editorial scope.} Counted unit: the article's theorem thm:meta-gaming (source0.tex lines 590--594). The article's complete original proof of it is delivered with it (source0.tex lines 605--701, one \texttt{proof} environment of 97 lines, which the article places away from the statement). No mathematics is written, repaired or re-derived: the statement and the proof are the article's own lines, in the article's own notation, and no retained line is substituted or re-flowed.  Retained uncounted as the source-local context the proof needs: lines 535--544; lines 547--556.  One declared adaptation: exactly 1 ancillary strings inside the retained spans are omitted (image-inclusion commands and, where the source repeats a label, the repeated label definitions) and no other line differs from the source; the figure environments, with their placement, captions and labels, are kept, and the package records the omitted strings in fidelity receipts bound to the affected bodies.  No retained line contains a citation, so the wrapper carries no bibliography and no bibliography entry.  The retained lines contain the article's own diagram code, which the wrapper typesets with the standard tikz packages; no image file is delivered and the package declares no asset dependency.  Editorial formatting only, in addition: an AMS \texttt{amsart} preamble that restates the article's own declarations and macros used by the retained lines, namely the environments \texttt{definition}, \texttt{lemma}, \texttt{theorem} on separate counters, together with the standard packages those lines need.  The retained environments print in this document as Definition 1, Lemma 1, Theorem 1 in order of appearance, whereas the article prints the same items with the numbering of its own full document.  The complete native source of the article as distributed in the arXiv e-print for this exact version is bundled unchanged as \texttt{sources/172887/source0.tex}.
\begin{definition}\label{def:quasi-homothety-seq}
    Let $A\geq 1, B\geq 0$, $\rho = (\rho_j)_j$ a sequence of positive integers with $\rho_j \to \infty$ and $\Gamma,\Delta_j$ graphs for $j\in \mathbb{N}$. A collection of $\iota_j \colon \Delta_j \to \Gamma$ is called a \emph{sequence of $(A\rho + B)$-surjective $(A,B)$-quasi-$\rho$-homotheties} if for all $j\in \mathbb{N}$
\begin{itemize}
    \item [(1)]  $\forall x \in \Gamma , \exists \overline{x} \in \Delta_j \colon
    d_{\Gamma}(\iota_j(\overline{x}),x) \leq A\rho_j + B$.

    \item [(2)] $\forall \overline{x}, \overline{y} \in \Delta_j \colon
    A^{-1}d_{\Delta_j}(\overline{x},\overline{y})- B \leq \frac{1}{\rho_j}d_{\Gamma}(\iota_j(\overline{x}), \iota_j(\overline{y})) \leq A \, d_{\Delta_j}(\overline{x},\overline{y}) + B $.
\end{itemize}
\end{definition}

\begin{lemma} \label{lem:quasi-homothety-seq}
    Let $\iota_j \colon \Delta_j \to \Gamma$ be a sequence of $(A\rho+B)$-surjective $(A,B)$-quasi-$\rho$-homotheties. Then there are maps $\pi_j \colon \Gamma \to \Delta_j $ such that for all $x,y \in \Gamma$
     $$
    \frac{1}{A} d_{\Delta_j}(\pi_j(x),\pi_j(y)) - (2A+3B) \leq \frac{1}{\rho_j} d_{\Gamma}(x,y) \leq A\, d_{\Delta_j}(\pi_j(x),\pi_j(y)) + (2A + 3B)
    $$
    and for all $x\in \Gamma , \overline{x} \in \Delta_j$
    $$
    d_{\Gamma}(x,\iota_j(\pi_j(x))) \leq A\rho_j+ B  \quad \text{ and } \quad d_{\Delta_j}(\overline{x}, \pi_j(\iota_j(\overline{x}))) \leq AB.
    $$
\end{lemma}

\begin{theorem}[Meta-gaming theorem]\label{thm:meta-gaming}
    Let $\Delta_j$ be graphs with $\wCop(\Delta_j) \geq m \in \mathbb{N}\cup \{\infty\}$ such that for every $\overline{\sigma}, \overline{\rho}$ there exist $\overline{\psi}, \overline{R}$ such that there exist winning strategies for the robber on $\Delta_j$ against $m$ cops, where if the cop picks speed \(\overline{\sigma}\) and reach \(\overline{\rho}\), the robber picks speed $\overline{\psi}$ and reach $\overline{R}$. 
    
    Let $A\geq 1, B\geq 0$ and $\rho_j \to \infty$ a sequence of positive integers. Let $\iota_j \colon \Delta_j \to \Gamma$ be a sequence of $(A\rho+B)$-surjective $(A,B)$-quasi-$\rho$-homotheties. Then $\sCop(\Gamma)  \geq m$.
\end{theorem}

\begin{proof}[Proof of \Cref{thm:meta-gaming}]


	Let \(n < m \leq \wCop(\Delta_j)\). For the constants $A,B$ in the statement, let $\overline{\sigma} := 4A^2+3AB$ and $\overline{\rho} = 4A(2A+3B)$. According to the assumptions there exist $\overline{\psi}$ and $\overline{R}$ for which the robber can win on $\Delta_j$ (for any $j$) against $n$ weak cops assuming that the cops started by picking $\overline{\sigma}$ and $\overline{\rho}$.

    We now describe a winning strategy for the robber player against $n$ strong cops on $\Gamma$. The game starts with the cop choosing some speed $\sigma \geq 1$. The robber chooses speed
    $$
    \psi = \sigma  \left( A \overline{\psi} + B\right) .
    $$
    The cop then chooses some reach $\rho$. Without loss of generality, we may assume that $\rho > \sigma$ and that $\rho = \rho_j$ for some $j$ large enough. The robber chooses the radius
    $$
    R = \rho (A\overline{R} + A + 2B).
    $$
    The cop then chooses starting positions $c_i \in \Gamma$ for $i = 1 , \ldots , n$. For the choice of $j$ with $\rho = \rho_j$ above, we fix $\Delta := \Delta_j$ and $\iota := \iota_j \colon \Delta \to \Gamma$. Let $\pi  \colon \Gamma \to \Delta$ be the quasi-projection from Lemma \ref{lem:quasi-homothety-seq}. Definition \ref{def:quasi-homothety-seq} and Lemma \ref{lem:quasi-homothety-seq} are summed up by the following facts for $x,y \in \Gamma, \overline{x}, \overline{y} \in \Delta$.
    \setcounter{equation}{0}
    \begin{align}
        &\forall x \in \Gamma \, \exists \overline{x} \in \Delta \colon d_{\Gamma}(\iota(\overline{x}),x) \leq A\rho + B \label{eq:def-surj} \\
    & \frac{1}{A} d_{\Delta}(\overline{x},\overline{y}) - B \leq \frac{1}{\rho}d_{\Gamma}(\iota(\overline{x}),\iota(\overline{y}))  \leq A \, d_{\Delta}(\overline{x},\overline{y}) + B \label{eq:def-quasi} \\
    &\frac{1}{A}  d_{\Delta}(\pi(x),\pi(y))  -  (2A + 3B) \leq \frac{1}{\rho} d_{\Gamma}(x,y)\leq A\,  d_{\Delta}(\pi(x),\pi(y))  +  (2A + 3B) \label{eq:lem-quasi} \\
    &    d_{\Gamma}(x,\iota(\pi(x))) \leq A\rho+ B \label{eq:lem-surj} \\
    & d_{\Delta}(\overline{x}, \pi(\iota(\overline{x})) ) \leq AB \label{eq:lem-proj}
    \end{align}
    During the game, the robber uses the weak-cop game played on $\Delta$ with parameters $\overline{\sigma}, \overline{\rho}, \overline{\psi} , \overline{R}$ as an oracle for how to play on $\Gamma$. To determine the robber's starting position $r_0$, the robber consults the oracle with cop starting positions $\pi(c_i) \in \Delta$ to obtain a starting position $\overline{r} \in \Delta$ with which the weak-cop game can be won on $\Delta$, in particular $d_\Delta(\overline{r}, \pi(c_i)) > \overline{\rho} + \overline{\sigma}$. The robber then chooses the actual starting position $r := \iota(\overline{r}) \in \Gamma$. By (\ref{eq:lem-quasi}) and (\ref{eq:lem-proj}),
    \begin{align*}
        d_{\Gamma}(\iota(\overline{r}),c_i) & \geq \rho \left( \frac{1}{A} \left( d_\Delta (\pi(\iota(\overline{r})), \pi(c_i) )- (2A + 3B) \right)\right) \\
        &\geq \rho \left( \frac{1}{A} ( d_\Delta(\overline{r}, \pi(c_i)) - d_\Delta(\overline{r} , \pi(\iota(\overline{r})) ) )-(2A+ 3B) \right) \\
        &> \rho \left( \frac{1}{A} (\overline{\rho} + \overline{\sigma} - AB ) -(2A + 3B) \right) \\
        &=  \rho\left( 4(2A + 3B) + \frac{1}{A}\overline{\sigma} - B - (2A + 3B) \right) \\
        & \geq \rho(2(2A + 3B)) \geq 2\rho \geq \rho + \sigma ,
    \end{align*}
    the robber is not caught at the starting position $r$, even after the cops move once.

    The game now proceeds in meta-stages. At the beginning of a meta-stage, the robber is placed at $r_t = \iota(\overline{r}_t)$ for some $\overline{r}_t \in \Delta$. The cops' positions $c_i$ are projected to $\Delta$, and the oracle suggests a path $ \overline{r}_t = \overline{p}_1 , \overline{p}_2, \ldots , \overline{p}_k \in \Delta$ along which the robber follows in a winning strategy on $\Delta$. The actual robber will now spend the next $\lambda := \lceil \rho/\sigma \rceil $ many turns following geodesic paths from $p_j := \iota(\overline{p}_j)$ to $p_{j+1}$ for $j \in \{ 1, \ldots , k-1\}$ to finally reach $p_k \in \iota(\Delta)$. During the meta-stage, the cops can also do $\lambda$ many moves. The meta-stage strategy then repeats. To prove that this gives a winning strategy for the robber it remains to show that
    \begin{itemize}
        \item [(a)] the robber can actually move as described from $p_1$ to $p_k$ in $\lambda$ many moves,
        \item [(b)] the movement of the cops for $\lambda$ many turns corresponds to a valid single move when projected to $\Delta$,
        \item [(c)] the robber is never caught during a meta-stage, and
        \item [(d)] after any number of turns, the robber eventually returns to the ball of radius $R$ around the base point.
    \end{itemize}

    (a) We have $d_\Delta(\overline{p}_1, \overline{p}_k) \leq \overline{\psi}$. Using (\ref{eq:def-quasi}) and $1 \leq \rho/\sigma \leq \lceil \rho / \sigma \rceil = \lambda$,
\begin{align*}
    d_{\Gamma}(p_1, p_k) & \leq \rho\left(A \, d_\Delta(\overline{p}_1, \overline{p}_k) + B \right)
    \leq \frac{\rho}{\sigma} \sigma \left(A \overline{\psi} + B\right)
    \leq \lambda \sigma (A \overline{\psi} + B) \leq \lambda \psi
\end{align*}
shows that the robber can complete the movement from $p_1$ to $p_k$ within $\lambda$ many turns.


\begin{figure}
  \centering
   % Replace 'example-image' with the filename of your image
  \caption{At the beginning of a meta-stage, the cop's position $c_0$ is provided to the oracle $\Delta$ via $\pi$. In the oracle there is a winning strategy by following along the path $\overline{p}_1, \ldots , \overline{p}_k$, which the robber implements in $\Gamma$, by following along $p_1 := \iota(\overline{p}_1), \ldots , p_k$. }
  \label{fig:meta_game}
\end{figure}
(b) If the starting position of a cop is $c_0\in \Gamma$, then their position $c_t \in \Gamma$ after $\lambda$ many turns satisfies $d_{\Gamma}(c_0,c_t) \leq \lambda \sigma $. Projected to $\Delta$, we use (\ref{eq:lem-quasi}) and $\lambda \sigma = \lceil \rho/\sigma \rceil \sigma \leq (\rho/\sigma + 1) \sigma = \rho + \sigma \leq 2\rho$ to show
\begin{align*}
    d_\Delta(\pi(c_0), \pi(c_t))
    & \leq A \left( \frac{1}{\rho} d_{\Gamma}(c_0,c_t) + 2A + 3B\right)
    \\
    &\leq A \left( 2 + 2A +3B \right) \leq A(4A +3B) = \overline{\sigma}
\end{align*}
which confirms that the projected movement of the cops corresponds to a valid movement in $\Delta$.



(c) Let $c_t$ be the position of a cop and $r_t$ the position of the robber at any time (possibly even different times) during a meta-stage. Note that in the projected game, the cops $\pi(c_t)$ are never too close to any of the vertices $\overline{p}_1, \ldots , \overline{p}_k$ of the path of the robber, namely
$$
d_\Delta(\pi(c_t), \overline{p}_i) > \overline{\rho}
$$
for every $i \in \{1, \ldots, k\}$. Since the robber in $\Gamma$ moves on geodesics along the $p_i$, there is a $j \in \{1, \ldots, k\}$ such that $d_{\Gamma}(r_t, p_j) \leq d_{\Gamma}(p_j, p_{j+1})$. The situation is illustrated in Figure \ref{fig:meta_game}. Using (\ref{eq:def-quasi}), we get an estimate
\begin{align*}
    d_{\Gamma}(r_t, p_j)
    &\leq d_{\Gamma}(\iota(\overline{p}_j), \iota(\overline{p}_{j+1}))
     \leq \rho ( A d_\Delta(\overline{p}_j, \overline{p}_{j+1}) + B ) = \rho(A + B)
\end{align*}
on the distance from the robber to some $p_j$. We then use the triangle inequality, the above estimates, (\ref{eq:def-quasi}), (\ref{eq:lem-surj}) and $\overline{\rho} = 4A(2A + 3B)$ to obtain
\begin{align*}
    d_{\Gamma}(r_t, c_t) & \geq d_{\Gamma}(\iota (\pi (c_t)) , \iota(\overline{p}_j)) - d_{\Gamma}(r_t , p_j) - d_{\Gamma}(\iota(\pi(c_t)), c_t) \\
    & \geq \rho \left(\frac{1}{A} d_\Delta(\pi(c_t), \overline{p}_j) - B \right) - \rho (A + B) - (A\rho + B) \\
    & > \rho \left( \frac{1}{A} \overline{\rho} - B \right) - 2\rho(A+B) \\
    & = \rho ( 4(2A+3B) -B -2A -2B) \geq 2\rho \geq \rho + \sigma
\end{align*}
which means that $r_t$ is always at least $\rho$ far away from $c_t$, even after one more movement by a distance of at most $\sigma$ by the cop. This means that the robber is never caught.

(d) Let $v$ be the base point in $\Gamma$ and $\pi(v)$ the base point in $\Delta$.
For any number of turns $k$, there is a number $k' \geq k$ such that the robber is in the ball of radius $\overline{R}$ around $\pi(v)$, that is $d_\Delta(\pi(v), \overline{r}) \leq \overline{R}$, where $\overline{r}$ is the position of the robber at the beginning of turn $k'$. In $\Gamma$, the robber will eventually visit $\iota(\overline{r})$, which lies in the ball by the following estimate. We use the triangle inequality, (\ref{eq:lem-surj}) and (\ref{eq:def-quasi}) to obtain
\begin{align*}
    d_{\Gamma}(v,\iota(\overline{r})) & \leq d_{\Gamma}(v, \iota(\pi(v))) + d_{\Gamma}( \iota(\pi(v)), \iota(\overline{r}))  \\
    &\leq ( A\rho + B) + \rho(A \, d_\Delta(\pi(v),\overline{r}) + B) \\
    &
    \leq  \rho (A + B + A\overline{R} + B)= R.
\end{align*}


The above calculations show that the robber can implement the strategy from the weak cop game on $\Delta$ using meta-stages to win the strong cop game on $\Gamma$. Hence $\operatorname{sCop}(\Gamma) > n$ for every $n< m = \wCop(\Delta)$, so $\sCop(\Gamma) \geq m$.
\end{proof}

\end{document}
